%% kapa: exam=1 round=19 year=2008 period=1 form=A subject=civil_law subject_ko=민법 seq=13 number=13
%% source: Q-Net 감정평가사 기출문제 (cst00309) artlSeq=1001962 file=extracted/1008/1001962/2008년 제19회 감정평가사 1차시험 기출문제/1교시.hwp
%% conversion: kordoc 4.24.51 -> tools/gen1.py (auto); review=proofread-sonnet; flags=none
\begin{question}{13}
원물과 과실에 관한 설명으로 옳은 것은? (다툼이 있는 경우에는 판례에 의함)
\begin{choices}
  \item 천연과실은 원물로부터 분리되기 전에는 원물의 구성부분에 지나지 않으므로 어떠한 경우에도 독립한 물건이 될 수 없다.
  \item 원물 사용의 대가를 받을 수 있는 권리도 과실이다.
  \item 천연과실은 수취할 권리의 존속기간 일수의 비율로 취득한다.
  \item 저당부동산에 관한 과실수취권은 목적물에 대한 압류가 있기 전에는 저당권설정자에게 있다.
  \item 물건의 사용이익은 과실이 아니므로 사용이익의 귀속에 관해서는 과실수취권에 관한 규정이 유추적용될 수 없다.
\end{choices}
\end{question}
\answer{4}
